\section{Visión del \texorpdfstring{\ac{GRID}}{GRID}}

La misión de la Universidad del Quindío se complementa con su visión institucional, la cual también es adoptada por el \ac{GRID}. A continuación se presenta la visión de la universidad:

\begin{quote}
	\textit{La Universidad del Quindío en el 2040 será una institución líder en procesos de impacto territorial, transformadora, innovadora, adaptable, con articulación y prospectiva glocal, que forme ciudadanos íntegros, comprometidos con la sociedad, con pensamiento crítico, sensibilidad social, y capacidad para afrontar los desafíos cambiantes del mundo. \citep{uniquindio2026}}
\end{quote}

A partir de esta visión, se identifican los siguientes enfoques estratégicos:

\begin{itemize}
	\item \textbf{Gestión:} Orientada hacia una institución líder, adaptable e innovadora, con capacidad para responder a los cambios y desafíos del entorno.

	\item \textbf{Docencia:} Enfocada en la formación de ciudadanos íntegros, con pensamiento crítico, sensibilidad social y capacidades para afrontar los desafíos cambiantes del mundo.

	\item \textbf{Investigación:} Orientada a generar conocimiento y capacidades que contribuyan a procesos de transformación e impacto territorial, articulándose con las necesidades y desafíos del entorno.

	\item \textbf{Extensión y Desarrollo Social:} Dirigida a fortalecer el impacto territorial de la Universidad mediante la articulación con su entorno y la generación de iniciativas que contribuyan a la transformación de la sociedad.

	\item \textbf{Responsabilidad Social:} Orientada a formar ciudadanos comprometidos con la sociedad, con sensibilidad frente a las necesidades de su entorno y capacidad para contribuir a procesos de transformación con impacto territorial.
\end{itemize}
